\chapter{Evaluation}
\label{chap:evaluation}

This chapter evaluates the executor module from a quantitative angle. The most
direct way to compare an ERC-7579 module against the way users currently
interact with AAVE is to measure gas: how much does a typical workflow cost
when issued through the module, and how does that compare with sending the
equivalent transactions from a regular externally owned account (EOA)? Rather
than building a dashboard that wraps the module in a user interface, I focus
on the cost of the underlying on-chain operations, because that cost is the
same regardless of which front-end is later layered on top.

The chapter first describes the methodology used to obtain the gas numbers,
then presents the results across four representative workflows. It compares
three configurations: an EOA calling AAVE directly, an ERC-7579 smart account
issuing the same calls as a native batched UserOperation, and the same smart
account issuing the workflow through my executor module. The chapter closes
with a comparison against related solutions in the AAVE tooling ecosystem and
a discussion of the limitations of the measurement.

\section{Methodology}

All measurements were produced by a Foundry test contract (\texttt{test/Gas.t.sol})
running against an Ethereum mainnet fork pinned at block 18,000,000. The fork
block, EVM version (Cancun) and RPC endpoint are fixed in \texttt{foundry.toml}
so the numbers are reproducible by re-running \texttt{forge test
--match-contract GasTest -vv}. The compiler is run without optimizer overrides
so the executor module is measured exactly as it is shipped in \texttt{src/}.

\subsection{What is measured}

For each workflow I record the total gas the user is responsible for paying.
This includes the intrinsic transaction cost of 21{,}000 gas per on-chain
transaction. The execution work (storage writes, oracle reads, AAVE bookkeeping)
is the same in all three configurations; the cost difference comes from how
many transactions are needed and from the overhead introduced by the
ERC-4337 EntryPoint when a smart account is involved.

For the EOA configuration, each AAVE interaction is a separate transaction. I
measure the pure execution gas of the calls with \texttt{gasleft()} snapshots
and then add 21{,}000 for each transaction the user would have to broadcast.
For the two smart account configurations, the entire workflow is bundled into a
single UserOperation routed through the EntryPoint via
\texttt{handleOps}, and the bundler transaction itself pays the 21{,}000
intrinsic cost once.

\subsection{Three configurations}

\textbf{EOA.} The user calls AAVE directly. Supplying requires two transactions
(\texttt{approve} followed by \texttt{Pool.supply}); borrowing requires one
transaction; repaying requires another approval plus a \texttt{Pool.repay}
transaction; withdrawing is a single transaction. This is the baseline that
every comparison against AAVE-tooling gas costs must use, because it is the
path most users take today.

\textbf{Native ERC-7579 batch.} The user owns an ERC-7579 smart account but
does \emph{not} install any application-specific module. They batch the same
raw approve and Pool calls into one UserOperation using the account's native
batched execution mode. This configuration isolates the cost of the account
abstraction framework itself, separate from any cost that my module adds on
top.

\textbf{Smart account with the executor module.} The same smart account,
with the executor module from \texttt{src/ExecutorTemplate.sol} installed.
Each user-level action (supply, borrow, repay, withdraw) is encoded as a
single \texttt{executor.execute} call, and multi-action workflows are bundled
as a batch of executor calls in one UserOperation. The module performs the
necessary approvals internally, so the user's calldata only carries the
high-level intent rather than the raw token-approval sequence.

\subsection{Workflows}

Four workflows of increasing length are measured:

\begin{enumerate}
    \item \textbf{Supply} — supply 10{,}000 USDC to the AAVE pool. The EOA
    path needs two transactions because USDC must first be approved.
    \item \textbf{Supply + borrow} — supply collateral, then borrow 0.5 WETH
    against it.
    \item \textbf{Supply + borrow + repay} — the above, followed by an
    approval of WETH and a repayment of the same 0.5 WETH.
    \item \textbf{Full cycle} — supply, borrow, repay, then withdraw half of
    the original collateral.
\end{enumerate}

These cover the full lifecycle implemented in the module and let me observe
how each configuration scales as more actions are added.

\section{Results}

Table~\ref{tab:gas-results} reports the total gas cost of each workflow under
each configuration. All numbers are in raw gas units; the bundler base fee is
already included for the smart account paths and the per-transaction base fee
is already included for the EOA path.

\begin{table}[ht]
\centering
\begin{tabular}{lrrr}
\hline
\textbf{Workflow} & \textbf{EOA} & \textbf{Native 7579} & \textbf{Module} \\
\hline
Supply                       &   251{,}824 &   502{,}911 &   525{,}475 \\
Supply + borrow              &   510{,}996 &   760{,}165 &   798{,}728 \\
Supply + borrow + repay      &   629{,}701 &   866{,}148 &   921{,}805 \\
Full cycle (supply, borrow, repay, withdraw) &   715{,}477 &   952{,}055 & 1{,}020{,}512 \\
\hline
\end{tabular}
\caption{Total gas cost per workflow, in gas units. EOA assumes one
transaction per AAVE interaction; smart account configurations bundle the
full workflow into a single UserOperation.}
\label{tab:gas-results}
\end{table}

\subsection{Smart account overhead is roughly constant}

The gap between the EOA and the smart account paths sits in a narrow band:
274{,}000 gas for the shortest workflow, 305{,}000 for the longest. This is
the overhead of running the same work through ERC-4337 rather than through a
plain transaction: signature validation, nonce checks, EntryPoint accounting
and the indirection between the EntryPoint, the account and the executor.
Because this overhead is paid once per UserOperation, it is amortised across
however many actions the workflow contains.

The break-even point follows directly from this. Each EOA transaction adds
21{,}000 gas in intrinsic cost. The smart account adds roughly 250{,}000 gas
in framework overhead. A workflow would need to contain more than
$\lceil 250{,}000 / 21{,}000 \rceil = 12$ transactions before the smart
account becomes cheaper than the EOA. None of the four AAVE workflows I
tested approach that threshold, and most realistic AAVE interactions a user
would issue in practice do not either.

\subsection{The module costs slightly more than a raw 7579 batch}

Comparing the third column with the second isolates the cost of the executor
module specifically. The module adds between 22{,}500 gas (single supply) and
68{,}500 gas (full cycle). This overhead comes from the additional contract
hop introduced by \texttt{executor.execute}: the account calls into the
module, the module decodes its action enum, and the module then calls back
into the account using \texttt{\_execute} to perform each underlying step.
Each round trip and each ABI decoding adds a small but measurable cost.

The cost is not material on a per-workflow basis -- under five percent for
the shortest workflow, under eight percent for the longest. What the user
gets in exchange is a smaller and safer surface to interact with: a single
function with a typed action enum, a pool address and asset addresses fixed
at install time, and approvals issued automatically rather than left for the
user to manage. Whether that trade-off is worth paying for is a UX question,
not a gas question.

\subsection{The smart account path is not a gas-saving optimisation}

The most important finding is that, for the workflows I measured, none of
the smart account configurations are cheaper than the corresponding EOA
workflow. This contradicts a common framing of account abstraction as a way
to reduce on-chain costs by batching. Batching does eliminate
per-transaction base fees, but only after enough operations have been
batched to outweigh the fixed UserOperation overhead. For a single
deposit-and-borrow that threshold is far out of reach.

The honest framing of this chapter, and of the project as a whole, is
therefore that the module is not justified by gas. It is justified by the
features that account abstraction enables and that the module exposes
cleanly to those features: paymaster sponsorship so a user can interact with
AAVE without holding ETH for fees, session keys so a delegated agent can
execute pre-approved AAVE actions on the user's behalf, and a typed action
interface that downstream tooling (validators, hooks, dashboards) can reason
about without parsing raw approve and Pool calldata.

\section{Comparison with other AAVE-tooling solutions}

The executor module is one of several patterns developed to make interacting
with AAVE more convenient than issuing raw transactions. This section places
it in context against three families of existing solutions. None of these
are benchmarked here in gas: they each rely on different account models and
deployment assumptions, and a fair head-to-head benchmark would require
re-deploying each of them at the same fork block. The comparison here is
qualitative.

\subsection{DeFi Saver and similar dashboards}

DeFi Saver and other dashboards built on top of AAVE do not change the
account model. The user still signs ordinary EOA transactions; the dashboard
simply assembles convenient routes (for example, leveraged-position
construction, automated liquidation protection) and prompts the user to
approve them. The result is excellent UX for a single user but still
requires multiple transactions and multiple signatures for any non-trivial
flow, and it does not give a third party the ability to act on the user's
behalf.

\subsection{Instadapp DSA}

Instadapp's Decentralised Smart Account (DSA) is closer in spirit to my
module: it deploys a smart contract per user that owns positions across
several DeFi protocols and exposes ``connector'' contracts which encapsulate
protocol-specific logic. A user can ask the DSA to execute a sequence of
connector calls in a single transaction. The architectural overlap with
ERC-7579 is significant -- both are smart-account-with-modules designs --
but the DSA predates the ERC-4337 ecosystem, lives behind a custom
permissioning model, and is not interoperable with bundlers or paymasters.
The trade-off the DSA makes is the opposite of what my module makes: it
gives up framework portability in exchange for being a self-contained
ecosystem.

\subsection{Safe modules}

Safe (formerly Gnosis Safe) is the most widely deployed smart-account
implementation, and modules can be installed on a Safe to extend its
behaviour. AAVE-specific Safe modules exist in the wild and are conceptually
similar to my executor: they receive a typed instruction and translate it
into the underlying approve and Pool calls. The difference is the framework.
Safe is not natively an ERC-4337 account and does not implement
ERC-7579's standardised module interface. A module written for Safe is not
portable to other accounts, whereas the module presented here implements
\texttt{ERC7579ExecutorBase} and runs unchanged on any ERC-7579-conforming
account (Kernel, Nexus, Safe7579, and so on). The portability is the main
contribution of building on top of ERC-7579 rather than on top of any single
smart-account implementation.

\section{Limitations}

The numbers reported in this chapter should be read with the following
caveats in mind.

\textbf{Single fork block.} All measurements were taken at block 18,000,000.
Gas costs for AAVE interactions can shift between blocks because of cold
versus warm storage access and because AAVE upgrades can change the
underlying logic. A campaign across many recent blocks would give a more
robust picture but would not change the qualitative shape of the comparison.

\textbf{Calldata cost.} The 21{,}000 base fee added to each EOA transaction
captures the intrinsic cost but not the per-byte calldata cost (4 gas per
zero byte, 16 gas per non-zero byte). For \texttt{approve} and the AAVE Pool
functions this adds roughly one to two thousand gas per transaction. The
EOA numbers in Table~\ref{tab:gas-results} are therefore a slight
under-count. Including calldata gas would make the EOA path marginally more
expensive but would not affect the overall conclusion.

\textbf{Bundler economics.} The smart account numbers report the gas the
EntryPoint charges. A real bundler typically marks up that figure with a
priority premium and a paymaster surcharge if one is in use. These mark-ups
are off-chain commercial choices and would not appear in any on-chain
benchmark.

\textbf{No paymaster, no session keys.} The benchmark exercises only the
plain user-signs-the-userOp path. The features that justify the smart
account approach -- gasless onboarding via paymasters, delegated signing
via session keys -- are not measured here because they shift the cost
question from ``how much gas?'' to ``who pays the gas, and what does the
user have to hold?''. That question is answered qualitatively in the
preceding section.

\textbf{Module configuration.} The executor stores the pool address and two
asset addresses at install time. The cost of installing the module is paid
once per smart account and is not amortised in the per-workflow numbers
above. For the test fixture, the install cost is roughly 25{,}700 gas
(from the \texttt{testOnInstall} measurement in
\texttt{test/ExecutorTemplate.t.sol}).

\section{Discussion}

The evaluation supports a single, honest conclusion: ERC-7579 plus an
application-specific executor module is not the right tool for users whose
goal is the lowest possible gas bill on a one-shot AAVE interaction. The
EOA path remains the cheapest way to issue any short workflow against AAVE
and is likely to remain so until the framework cost of an ERC-4337
UserOperation drops by an order of magnitude.

The tool is, however, the right shape for a different goal: giving users
and the agents that act on their behalf a single, typed handle on AAVE that
plays well with the rest of the account-abstraction stack. The module
offers a stable interface (\texttt{execute(action, ...)}), encapsulates the
correctness-critical sequencing of approve and Pool calls, and inherits all
of the properties of the surrounding ERC-7579 ecosystem -- programmable
validators, hooks, session keys, and paymasters -- without writing a single
custom line for any of them. The 5--8\% overhead the module adds on top of
a raw batch is a fair price for that surface, and the 30--40\% overhead
the smart account path adds on top of the EOA path is the price one pays
for the features themselves, not for the module.
